%! Choosing and Comparing Models in Context — Unit 2, Lesson 2.7 — Exit Ticket (Answer Key)
\documentclass[10pt]{article}
\usepackage{atda-article}
\usepackage{atda-key}   % pulls in -boxes; do NOT also load -boxes
\newcommand{\TallMath}[1]{$\displaystyle #1\rule[-1.4em]{0pt}{3.2em}$}
\setlength{\workrowsep}{8pt}

\begin{document}

\pageheader{Unit 2, Lesson 2.7}{Exit Ticket --- Answer Key}
% NAMESTRIP: no \namedateperiod here — mirrors the blank. Teacher-only prose goes
% in the lesson plan as \begin{teachernote}[Exit Ticket], never in a key.

\vspace{0.15in}

\begin{notesbox}{One Gym, Two Tests, and a Building That Holds $500$}
\small
Notes closed. Three items. A gym opened five months ago and counted its members at the end of each
month.

\vspace{4pt}
\renewcommand{\arraystretch}{1.3}
\begin{center}
\begin{tabular}{@{}|l|c|c|c|c|c|@{}}
\hline
Month $m$ & $0$ & $1$ & $2$ & $3$ & $4$ \\ \hline
Members & $64$ & $96$ & $144$ & $216$ & $324$ \\ \hline
First differences & --- & \ans{$+32$} & \ans{$+48$} & \ans{$+72$} & \ans{$+108$} \\ \hline
Second differences & --- & --- & \ans{$+16$} & \ans{$+24$} & \ans{$+36$} \\ \hline
Ratios & --- & \ans{$1.5$} & \ans{$1.5$} & \ans{$1.5$} & \ans{$1.5$} \\
\hline
\end{tabular}
\end{center}

\begin{enumerate}[leftmargin=1.6em, itemsep=7pt, topsep=6pt]

\item Fill the three rows above. Then name the family that models this gym, and name the row that
decided it.

\ansline{\textbf{Exponential.} The \textbf{ratios} row is the constant one --- every month is $1.5$ times the month before,}\\
\ansline{which is a $50\%$ gain each month. The first and second differences both keep growing.}\\

\item Use your model to predict the membership at the end of month $5$ and month $6$. The building
is licensed to hold $500$ members. According to the model, which month is the first one over the
limit --- and what does that tell you about how far ahead this model can be trusted?

\ansline{Month $5$: $324 \times 1.5 = 486$ members. Month $6$: $486 \times 1.5 = 729$ members.}\\
\ansline{\textbf{Month 6} is the first one over $500$ --- and the building cannot hold $729$ people.}\\
\ansline{So the model is already predicting something impossible two months past the data: growth of $50\%$ a month has to stop.}\\

\item A classmate writes: \emph{``The gym gains more members every month --- $32$, then $48$, then
$72$. Growth that speeds up is quadratic.''} Noticing that the jumps get bigger was correct in item
1. Explain what is wrong with the conclusion anyway, and write the correct statement.

\ansline{Growing jumps only rule out \emph{linear}. Quadratic and exponential both have growing first differences,}\\
\ansline{so that observation cannot choose between them. The second differences here are $16, 24, 36$ --- not constant.}\\
\ansline{Correct: the \textbf{ratios} are constant at $1.5$, so the model is \textbf{exponential}, not quadratic.}\\

\end{enumerate}
\end{notesbox}

\end{document}
